


\subsection{Proof of Lemma \ref{lem:AO}}\label{sec:proofAO}


\vp
%\noindent\underline{Proof of (i):}~

\subsubsection{Proof of (i).} Strong convexity of the objective function in \eqref{eq:AO_con} is easily verified by the second derivative test and use of Assumption \ref{ass:mu} that $\lai\leq\Sigma_{\max},~i\in[p].$ Uniqueness of the solution follows directly from strong convexity. \ct{Strictly speaking we might need to also argue existence, i.e., feasibility of the AO. An indirect way is to show feasibility using the CGMT, but it  seems unnecessarily complicated?}



\vspace{5pt}
\subsubsection{Proof of (ii).}Using Lagrangian formulation, the solution $\wha_n$ to \eqref{eq:AO_con} is the same as the solution to the following:
\begin{align}
\left(\wha_n,u_n\right) :=\arg\min_{\w\in\Bc}\max_{u\geq 0} ~\frac{1}{2}\tn{\Sigmab^{-1/2}\w+\betas}^2 + u \left( \sqrt{\tn{\w}^2+\sigma^2} \tn{\gba} - \sqrt{\kappa}\,\hba^T\w + \frac{\sigma h}{\sqrt{n}} \right)\label{eq:AO_2}
\end{align}
where we have: (i) set $\gba := \g/\sqrt{n}$ and $\hba:= \h/\sqrt{p}$; (ii) recalled that $p/n=\kappa$; and, (iii) used $\left(\wha_n,u_n\right)$ to denote the optimal solutions in \eqref{eq:AO_2}. The subscript $n$ emphasizes the dependence of $\left(\wha_n,u_n\right)$ on the problem dimensions. Also note that (even though not explicit in the notation) $\left(\wha_n,u_n\right)$ are random variables depending on the realizations of $\gba,\hba$ and $h$.

 Notice that the objective function above is convex in $\w$ and linear (thus, concave) in $u$. Also, $\Bc$ is compact. Thus, strong duality holds and we can flip the order of min-max \cite{fan1953minimax}. Moreover, in order to make the objective easy to optimize with respect to $\w$, we use the following variational expression for the square-root term $\sqrt{\tn{\w}^2+\sigma^2}$:
$$
\tn{\gba}\sqrt{\tn{\w}^2+\sigma^2} = \tn{\gba}\cdot\min_{\tau\in[\sigma,\sqrt{\sigma^2+B_+^2}]} \left\{ \frac{\tau}{2} + \frac{\tn{\w}^2+\sigma^2}{2\tau} \right\} = \min_{\tau\in[\sigma,\sqrt{\sigma^2+B_+^2}]} \left\{ \frac{\tau\tn{\gba}^2}{2} + \frac{\sigma^2}{2\tau} + \frac{\tn{\w}^2}{2\tau} \right\},
$$
where $B_+$ is defined in Lemma \ref{lem:bd_PO}. For convenience define the constraint set for the variable $\tau$ as $\Tc':=[\sigma,\sqrt{\sigma^2+B_+^2}]$. For reasons to be made clear later in the proof (see proof of statement (iii)), we consider the (possibly larger) set:
\[
\Tc:=[\sigma,\max\{\sqrt{\sigma^2+B_+^2},2\tau_*\}]\,
\]
where $\tau_*$ is as in the statement of the lemma.

The above lead to the following equivalent formulation of \eqref{eq:AO_2}, 
\begin{align}
\left(\wha_n,u_n,\tau_n\right) = \max_{u\geq 0}\min_{\w\in\Bc,\tau\in\Tc} ~ \frac{u\tau\tn{\gba}^2}{2} + \frac{u\sigma^2}{2\tau} + \frac{u\sigma h}{\sqrt{n}} + \min_{\w\in\Bc} \left\{ \frac{1}{2}\tn{\Sigmab^{-1/2}\w+\betas}^2 + \frac{u}{2\tau}\tn{\w}^2 - u \sqrt{\kappa}\,{\hba}^T\w \right\} .\label{eq:AO_3}
\end{align}
%\ct{To be fully rigorous, need to show here that the unconstrained min over $\w$ is the same as the constrained $\w\in\Bc$.}
The minimization over $\w$ is easy as it involves a strongly convex quadratic function. First, note that the unconstrained optimal $\w':=\w'(\tau,u)$ (for fixed $(\tau,u)$) is given by
\begin{align}\label{eq:w'}
\w':=\w'(\tau,u) = -\left(\Sigmab^{-1}+\frac{u}{\tau}\Iden\right)^{-1}\left(\Sigmab^{-1/2}\betas-u\sqrt{\kappa}\hba\right), 
\end{align}
and \eqref{eq:AO_3} simplifies to
\begin{align}
\left(u_n,\tau_n\right)=\max_{u\geq 0}\min_{\tau\in\Tc} ~ \frac{u\tau\tn{\gba}^2}{2} + \frac{u\sigma^2}{2\tau} + \frac{u\sigma h}{\sqrt{n}} - \frac{1}{2} \left(\Sigmab^{-1/2}\betas-u\sqrt{\kappa}\hba\right)^T\left(\Sigmab^{-1}+\frac{u}{\tau}\Iden\right)^{-1} \left(\Sigmab^{-1/2}\betas-u\sqrt{\kappa}\hba\right)\,=:\Rc(u,\tau) .\label{eq:AO_4}
\end{align}
%\somm{Does $u=0$ scenario create a problem in saddle point uniqueness?} 
It can be checked by direct differentiation and the second-derivative test that the objective function in \eqref{eq:AO_4} is strictly convex in $\tau$ and strictly concave in $u$ over the domain $\{(u,\tau)\in\R_+\times\R_+\}$ \footnote{{To analyze the matrix-vector product term in \eqref{eq:AO_4} for $(\tau,u)$ one can use the fact that $\bSi$ is diagonal. This way, as a function of $u$ and $\tau$ the analysis reduces to the properties of relatively simple functions. For instance, for $\tau$, this function is in the form $f(\tau)=-(a+b/\tau)^{-1}$ for $a,b>0$, which is strictly convex.}}. Thus, the saddle point $(u_n,\tau_n)$ is unique. Specifically, this implies that the optimal $\wha_n$ in \eqref{eq:AO_3} is given by (cf. \eqref{eq:w'})
\begin{align}\label{eq:w_n}
\wha_n=\w'(\tau_n,u_n) = -\left(\Sigmab^{-1}+\frac{u_n}{\tau_n}\Iden\right)^{-1}\left(\Sigmab^{-1/2}\betas-u_n\sqrt{\kappa}\hba\right).
\end{align}
In Lemma \ref{lem:AO}(v) we will prove that wpa 1, in the limit of $p\rightarrow\infty$, $\|\wha_n\|_2\leq C$ for sufficiently large absolute constant $C>0$. Thus, by choosing the upper bound in the definition of $\Bc$ in Lemma \ref{lem:bd_PO} strictly larger than C, guarantees that the unconstrained $\wha_n$ in \eqref{eq:w_n} is feasible in \eqref{eq:AO_3}.

\noindent{\textbf{Asymptotic limit of the key quantities $\tau_n,u_n$:}} In what follows, we characterize the high-dimensional limit of the optimal pair $(u_n,\tau_n)$ in the limit $n,p\rightarrow\infty,~p/n\rightarrow\kappa$.
%Now that we have reached a min-max problem over only scalar variables, we are ready to study its convergence properties in the limit $n,p\rightarrow\infty,~p/n\rightarrow\kappa$. 
We start by analyzing the (point-wise) convergence of $\Rc(u,\tau)$. 
For the first three summands in \eqref{eq:AO_4}, we easily find that
$$
\left\{\frac{u\tau\tn{\gba}^2}{2} + \frac{u\sigma^2}{2\tau} + \frac{u\sigma h}{\sqrt{n}}\right\}~ \rP~\left\{
\frac{u\tau}{2} + \frac{u\sigma^2}{2\tau} \right\}.
$$
Next, we study the fourth summand. First, note that
\begin{align}
(u\sqrt{\kappa}\hba)^T\left(\Sigmab^{-1}+\frac{u}{\tau}\Iden\right)^{-1}(u\sqrt{\kappa}\hba) &= u^2\kappa\,\frac{1}{p}\h^T\left(\Sigmab^{-1}+\frac{u}{\tau}\Iden\right)^{-1}\h \nn\\
&= u^2\kappa\,\frac{1}{p}\sum_{i=1}^{p}\frac{\h_i^2}{\lai^{-1}+\frac{u}{\tau}} \nn\\
&\rP u^2\kappa\,\E\left[ \frac{1}{\Lambda^{-1}+\frac{u}{\tau}} \right].
\end{align}
% the quadratic form
%\begin{align}
%(u\sqrt{\kappa}\hba)^T\left(\Sigmab^{-1}+\frac{u}{\tau}\Iden\right)^{-1}(u\sqrt{\kappa}\hba) = u^2\kappa\,\frac{1}{p}\h^T\left(\Sigmab^{-1}+\frac{u}{\tau}\Iden\right)^{-1}\h
%\end{align}
%concentrates around its expectation \cite{?} $u^2\kappa\,\tr\left(\left(\Sigmab^{-1}+\frac{u}{\tau}\Iden\right)\right)^{-1}.$

%\somm{Explain/justify PL(2)'ness of these functions. Clearly explain how A2 is used\dots}
In the last line, $\Lambda$ is a random variable as in Definition \ref{def:Xi}. {Also, we used Assumption \ref{ass:mu} together with the facts that $\h$ is independent of $\Sigmab$ and  that the function $(x_1,x_2)\mapsto x_1^2(x_2^{-1}+u/\tau)^{-1}$ is $\rm{PL}(3)$ assuming $x_2$ is bounded (see Lemma \ref{lem:PLbdd} for proof).}
\ct{Question: Is it immediate that the empirical distribution of $(\beta,\Sigmab,\h)$ converges in $W_k$ to $\mu\otimes\Nn(0,1)$ given that $(\beta,\Sigmab)$ converges to $\mu$ and $\h$ is independent???}
%
 Second, we find that
\begin{align}
(\betas)^T\Sigmab^{-1/2}\left(\Sigmab^{-1}+\frac{u}{\tau}\Iden\right)^{-1}\Sigmab^{-1/2}\betas &= \frac{1}{p}(\sqrt{p}\betas)^T\left(\Iden+\frac{u}{\tau}\Sigmab\right)^{-1}(\sqrt{p}\betas) \nn\\
&= \frac{1}{p}\sum_{i=1}^{p}\frac{\left(\sqrt{p}\betas_i/\sqrt{\lai}\right)^2}{\lai^{-1}+\frac{u}{\tau}}\nn\\
&\rP \E\left[\frac{B^2\Lambda^{-1}}{\Lambda^{-1}+\frac{u}{\tau}}\right].
\end{align}
Here, $\Lambda,B$ are random variables as in Definition \ref{def:Xi} and {we also used Assumption \ref{ass:mu} together with the fact that the function $(x_1,x_2)\mapsto x_1^2x_2^{-1}(x_2^{-1}+u/\tau)^{-1}$ is $\rm{PL}(3)$ assuming $x_2$ is bounded (see Lemma \ref{lem:PLbdd} for proof).}
Third, \cts{by independence of $(\betas, \Sigmab)$ from $\h$}
\begin{align}
(u\sqrt{\kappa}\hba)^T\left(\Sigmab^{-1}+\frac{u}{\tau}\Iden\right)^{-1}\Sigmab^{-1/2}\betas = u\sqrt{\kappa} \cdot \frac{1}{p}\sum_{i=1}
^p \frac{\h_i\Sigmab_{i,i}^{-1/2}(\sqrt{p}\betas_i)}{\Sigmab_{i,i}^{-1}+\frac{u}{\tau}\Iden} ~\rP~ 0.
\end{align}
Putting these together, the objective $\Rc(u,\tau)$ in \eqref{eq:AO_4} converges point-wise in $u,\tau$ to
\begin{align}
\Rc(u,\tau)\rP\Dc(u,\tau) := \frac{1}{2}\left({u\tau} + \frac{u\sigma^2}{\tau} - u^2\kappa\,\E\left[ \frac{1}{\Lambda^{-1}+\frac{u}{\tau}} \right] - \E\left[\frac{B^2\Lambda^{-1}}{\Lambda^{-1}+\frac{u}{\tau}}\right] \right).\label{eq:conv_pt}
\end{align}
Note that $\Rc(u,\tau)$ (and thus, $\Dc(u,\tau)$) is convex in $\tau$ and concave in $u$. Thus, the convergence in \eqref{eq:conv_pt} is in fact uniform (e.g., \cite{AG1982}) and we can conclude that 
\begin{align}\label{eq:Dc0}
\phi(\g,\h) \rP \max_{u\geq 0}\min_{\tau\in\Tc}~ \Dc(u,\tau).
\end{align}
and {using strict concave/convexity of $\Dc(u,\tau)$, we also have the parameter convergence \cite[Lem. 7.75]{NF36}}
%\somm{This line follows from strict convex/concavity right?}
\begin{align}\label{eq:Dc}
{(u_n,\tau_n) \rP (u_*,\tau_*):=\arg\max_{u\geq 0}\min_{\tau\in\Tc}~ \Dc(u,\tau).}
\end{align}
In the proof of statement (iii) below, we show that the saddle point of \eqref{eq:Dc0} is $(u_*,\tau_*)$. In particular, $\tau_*$ is strictly in the interior of $\Tc$, which combined with convexity implies that
$$
 \max_{u\geq 0}\min_{\tau\in\Tc}~ \Dc(u,\tau) =  \max_{u\geq 0}\min_{\tau>0}~ \Dc(u,\tau) =: \bar\phi.
$$
This, together with the first display above proves the second statement of the lemma.

\vspace{5pt}
\subsubsection{Proof of (iii).}  Next, we compute the saddle point $(u_*,\tau_*)$ by studying the first-order optimality conditions of the strictly concave-convex $\Dc(u,\tau)$. Specifically, we consider the unconstrained minimization over $\tau$ and we will show that the minimum is achieved in the strict interior of $\Tc$. Direct differentiation of $\Dc(u,\tau)$ gives
\begin{subequations}
\begin{align}
{\tau} + \frac{\sigma^2}{\tau} - 2u\kappa\E\left[ \frac{1}{\Lambda^{-1}+\frac{u}{\tau}} \right] + \frac{u^2}{\tau}\kappa \E\left[ \frac{1}{\left(\Lambda^{-1}+\frac{u}{\tau}\right)^2}\right] + \frac{1}{\tau}\E\left[\frac{B^2\Lambda^{-1}}{\left(\Lambda^{-1}+\frac{u}{\tau}\right)^2}\right] &= 0, \label{eq:fo1}\\
{u} - \frac{u\sigma^2}{\tau^2} - \frac{u^3}{\tau^2}\kappa \E\left[ \frac{1}{\left(\Lambda^{-1} + \frac{u}{\tau}\right)^2}\right] - \frac{u}{\tau^2} \E\left[\frac{B^2\Lambda^{-1}}{\left(\Lambda^{-1}+\frac{u}{\tau}\right)^2}\right] &= 0,\label{eq:fo2}
\end{align}
\end{subequations}
Multiplying \eqref{eq:fo2}  with $\frac{\tau}{u}$ and adding to \eqref{eq:fo1} results in the following equation
\begin{align}
\tau = u\kappa\E\left[ \frac{1}{\Lambda^{-1}+\frac{u}{\tau}} \right] ~\Leftrightarrow ~
\E\left[ \frac{1}{(\frac{u}{\tau}\Lambda)^{-1}+1} \right] = \frac{1}{\kappa} \label{eq:tauu}\,.
\end{align}
Thus, we have found that the ratio $\frac{u_*}{\tau_*}$ is the unique solution to the equation in \eqref{eq:tauu}. Note that this coincides with the Equation \eqref{eq:ksi} that defines the parameter $\xi$ in Definition \ref{def:Xi}.  The fact that \eqref{eq:tauu} has a unique solution for all $\kappa>1$ can be easily seen as $F(x)=\E\left[ \frac{1}{(x\Lambda)^{-1}+1} \right], x\in\R_+$ has range $(0,1)$ and is strictly increasing (by differentiation). 


Thus, we call $\xi=\frac{u_*}{\tau_*}$. Moreover, multiplying \eqref{eq:fo2} with $u$ leads to the following equation for $\tau_*$:
\begin{align}
u_*^2 = \sigma^2\xi^2 + u_*^2 \xi^2  \kappa \E\left[ \frac{1}{\left(\Lambda^{-1} +\xi\right)^2}\right] + \xi^2  \E\left[ \frac{B^2\Lambda^{-1}}{\left(\Lambda^{-1} + \xi\right)^2}\right] ~\Rightarrow~ \tau_*^2 = \frac{\sigma^2 + \E\left[ \frac{B^2\Lambda^{-1}}{\left(\Lambda^{-1} + \xi\right)^2}\right]}{1-\xi^2\kappa \E\left[ \frac{1}{\left(\Lambda^{-1} +\xi\right)^2}\right]} =  \frac{\sigma^2 + \E\left[ \frac{B^2\Lambda^{-1}}{\left(\Lambda^{-1} + \xi\right)^2}\right]}{1-\kappa \E\left[ \frac{1}{\left((\xi\Lambda)^{-1} +1\right)^2}\right]}.
\end{align}
{Again, note that this coincides with Equation \eqref{eq:gamma} that determines the parameter $\gamma$ in Definition \ref{def:Xi}, i.e., $\tau_*^2 = \gamma.$ By definition of $\Tc$ and of $\tau_*$, it is clear that $\tau_\star$ is in the strict interior of $\Tc$.}

\vspace{5pt}
\subsubsection{Proof of (iv).}  For convenience, define 
$$F_n(\btha,\betas,\Sigmab):=\frac{1}{p} \sum_{i=1}^{p} f\left(\sqrt{p}\btha_{n,i},\sqrt{p}\betas_{i},\Sigmab_{ii}\right)\quad\text{and}\quad\alpha_*:=\E_\mu\left[f(X_{\kappa,\sigma^2}(\Lambda,B,H),B,\Lambda)\right].$$
Recall from \eqref{eq:w_n} the explicit expression for $\wha_n$, repeated here for convenience.
\begin{align}
\wha_n = -\left(\Sigmab^{-1}+\frac{u_n}{\tau_n}\Iden\right)^{-1}\left(\Sigmab^{-1/2}\betas-u_n\sqrt{\kappa}\hba\right).\nn
\end{align}
Also, recall that $\btha_n = \Sigmab^{-1/2}\wha_n+\betas$. Thus,  (and using the fact that $\hba$ is distributed as $-\hba$),
\begin{align}
\btha_n &=\Sigmab^{-1/2}\left(\Sigmab^{-1}+\frac{u_n}{\tau_n}\Iden\right)^{-1}u_n\sqrt{\kappa}\hba + \left(\Iden-\Sigmab^{-1/2}\left(\Sigmab^{-1}+\frac{u_n}{\tau_n}\Iden\right)^{-1}\Sigmab^{-1/2}\right)\betas\nn\\
\Longrightarrow\btha_{n,i}&=\frac{{\Sigmab_{i,i}^{-1/2}}}{1+(\ksi_n\Sigmab_{i,i})^{-1}}\sqrt{\kappa}\tau_n\hba_i +\left(1-\frac{1}{1+\ksi_n\Sigmab_{i,i}}\right)\betas_i.\label{eq:betan}
\end{align}
For $i\in[p]$, define 
\begin{align}
\vb_{n,i} = \frac{{\Sigmab_{i,i}^{-1/2}}}{1+(\ksi_*\Sigmab_{i,i})^{-1}}\sqrt{\kappa}\tau_*\hba_i +\left(1-\frac{1}{1+\ksi_*\Sigmab_{i,i}}\right)\betas_i
%-\left(\Sigmab^{-1}+\frac{u_*}{\tau_*}\Iden\right)^{-1}\left(\Sigmab^{-1/2}\betas-u_*\sqrt{\kappa}\hba\right).
\end{align}
In the above, for convenience, we have denoted $\xi_n:=u_n/\tau_n$ and recall that $\xi_*:=u_*/\tau_*$. 


The proof proceeds in two steps. In the first step, we use the fact that $\ksi_n\rP\ksi_*$ and $u_n\rP u_\star$ (see \eqref{eq:Dc}) to prove that for any $\eps\in(0,\ksi_*/2)$, there exists an absolute constant $C>0$ such that wpa 1:
\begin{align}\label{eq:ivstep1}
|F_n(\sqrt{p}\btha_n,\sqrt{p}\betas,\Sigmab) - F_n(\sqrt{p}\vb_n,\sqrt{p}\betas,\Sigmab) | \leq C\eps.
\end{align}
In the second step, we use pseudo-Lipschitzness of $f$ and Assumption \ref{ass:mu} to prove that
\begin{align}
|F_n(\vb_n,\betas,\Sigmab)| \rP \alpha_*.\label{eq:ivstep2}
\end{align}
The desired follows by combining \eqref{eq:ivstep1} and \eqref{eq:ivstep2}. Thus, in what follows, we prove \eqref{eq:ivstep1} and \eqref{eq:ivstep2}.


\vspace{2pt}
\noindent\underline{Proof of \eqref{eq:ivstep1}.}~~Fix some $\eps\in(0,\ksi_*/2)$. From \eqref{eq:Dc}, we know that w.p.a. 1 $|\xi_n-\xi_*|\leq \eps$ and $|u_n-u_*|\leq \eps$. Thus, $\wha_n$ is close to $\vb_n$. Specifically, in this event, for every $i\in[p]$, it holds that:
\begin{align}
\nn|\btha_{n,i} - \vb_{n,i}| &\leq {|\betas_i|}\left|\frac{1}{1+\xi_n\lai}-\frac{1}{1+\xi_*\lai}\right| + \sqrt{\kappa}\frac{|\hba_i|}{\sqrt{\Sigmab_{i,i}}}\left|\frac{\tau_n}{1+(\xi_n\lai)^{-1}}-\frac{\tau_*}{1+(\xi_*\lai)^{-1}}\right| \\
\nn&= {|\betas_i|}\left|\frac{1}{1+\xi_n\lai}-\frac{1}{1+\xi_*\lai}\right| + \sqrt{\kappa}\frac{|\hba_i|}{\sqrt{\Sigmab_{i,i}}}\left|\frac{u_n}{\xi_n+\lai^{-1}}-\frac{u_*}{\xi_*+\lai^{-1}}\right| \\
\nn& \leq
{|\betas_i|}\frac{|\lai||\xi_n-\xi_*|}{|1+\xi_n\lai||1+\xi_*\lai|} + \sqrt{\kappa}\frac{|\hba_i|}{\sqrt{\lai}}  \frac{u_*|\ksi_n-\ksi_*|}{(\ksi_n+\lai^{-1})(\ksi_*+\lai^{-1})} + 
\sqrt{\kappa}\frac{|\hba_i|}{\sqrt{\lai}}\frac{|u_n-u_*|}{\xi_n+\lai^{-1}}
\\
\nn& \leq
{|\betas_i|}{\Sigma_{\max}\eps} + \sqrt{\kappa}{|\hba_i|} u_*\Sigma_{\max}^{3/2} \eps+ 
\sqrt{\kappa}|\hba_i|\Sigma_{\max}^{1/2}\eps
\\
&\leq C\eps  \left(|\hba_i| + |\betas_i|\right).\label{eq:betav}
%\cdot \max\left\{\Sigma_{\max}^{3/2},\Sigma_{\max}^{1/2}\right\}\,
%\nn& \leq
%\frac{|\betas_i|}{\sqrt{\lai}} \frac{\eps}{(\xi-\eps)\xi} + \sqrt{\kappa}|\hba_i|\left(\frac{\eps(u_*+\eps)}{(\xi-\eps)\xi} + \frac{\eps}{\xi}\right)\\
%& \leq
%\frac{|\betas_i|}{\sqrt{\lai}} c_1(\eps) + |\hba_i|c_2(\eps).\label{eq:wivi}
\end{align}%\som{$\eps$ missing}
where $C=C(\Sigma_{\max},\kappa,u_*)$ is an absolute constant. In the second line above, we recalled that $u_n=\tau_n\xi_n$ and $u_*=\tau_*\xi_*$. In the third line, we used the triangle inequality. In the fourth line, we used that $\ksi_*>0$, $0<\lai\leq\Sigma_{\max}$ and $\ksi_n\geq \ksi_*-\eps \geq \ksi_*/2 >0$. 
%In the second line above, we used the triangle inequality. In the third inequality we used Assumption \ref{ass:inv} that $\lai>0, i\in[p]$, and in the last line we defined appropriate constants $c_1(\eps),c_2(\eps)>0$. Importantly,  ... 

Now, we will use this and Lipschitzness of $f$ to argue that there exists absolute constant $C>0$ such that wpa 1,
\begin{align}
|F_n(\sqrt{p}\btha_n,\sqrt{p}\betas,\Sigmab) - F_n(\sqrt{p}\vb_n,\sqrt{p}\betas,\Sigmab) | \leq C\eps.\nn
\end{align}

Denote, $\ab_i=(\sqrt{p}\btha_{n,i},\sqrt{p}\betas_i,\Sigmab_{i,i})$ and $\bb_i=(\sqrt{p}\vb_{n,i},\sqrt{p}\betas_i,\Sigmab_{i,i})$. Following the exact same argument as in \eqref{eq:dev2show} (just substitute $\bt\leftrightarrow\vb_n$ in the derivation), we have that for some absolute constant $C>0$ wpa 1:
\begin{align}
|F_n(\btha_n(\g,\h),\betas,\Sigmab) - F_n(\vb_n,\betas,\Sigmab)| &\leq C \|\btha_n-\vb_n\|_2.
\end{align}
%\begin{align}
%|F_n(\btha_n(\g,\h),\betas,\Sigmab) - F_n(\vb_n,\betas,\Sigmab)| &\leq \frac{L}{p}\sum_{i=1}^p (1+\max\{\|\ab_i\|_2^{k-1},\|\bb_i\|_2^{k-1}\})\|\ab_i-\bb_i\|_2\nn \\
%&\leq \frac{L}{p}\sum_{i=1}^p \left(1+\max\{\|\ab_i\|_2^{k-1},\|\bb_i\|_2^{k-1}\}\right)\cdot|\sqrt{p}\btha_{n,i}(\g,\h) - \sqrt{p}\vb_{n,i}|\nn\\
%&\leq L\left(1+\max\{\frac{1}{p}\sum_{i=1}^p \|\ab_i\|_2^{2k-2},\frac{1}{p}\sum_{i=1}^p\|\bb_i\|_2^{2k-2}\} \right)^{1/2}{\|\btha_n(\g,\h) - \vb_n\|_2}\nn.
%\end{align}
%\max\left\{\Sigma_{\max}^{3/2},\Sigma_{\max}^{1/2}\right\}
From this and \eqref{eq:betav}, we find that
\begin{align}
|F_n(\btha_n(\g,\h),\betas,\Sigmab) - F_n(\vb_n,\betas,\Sigmab)|
&\leq  
C\eps \left(\sum_{i=1}^p \left(|\hba_i|+|\betas_i|\right)^2\right)^{1/2}\nn
\\
 &\leq 
C\eps\sqrt{2}\sqrt{\|\betas\|_2^2 + \|\hba\|_2^2}.\label{eq:epsS}
\end{align}
%As in \eqref{eq:bdmom}, it can be shown that wpa 1: $\frac{1}{p}\sum_{i=1}^p \|\ab_i\|_2^{2k-2}<\infty$ and $\frac{1}{p}\sum_{i=1}^p \|\bb_i\|_2^{2k-2}<\infty$, as $p\rightarrow\infty$. 
But, \cts{recall that $\|\betas\|_2^2 = \frac{1}{p}\sum_{i=1}^p(\sqrt{p}\betas_i)^2<\infty$, as $p\rightarrow \infty$ by Assumption \ref{ass:mu}.}
%{\color{red} by assumption on second moment convergence of $\sqrt{p}\betas$.} 
Also, since $\hba_i\sim\Nn(0,1/p)$, it holds that $\|\hba\|_2^2\leq 2$, wpa 1 as $p\rightarrow\infty$. Therefore, from \eqref{eq:epsS}, wpa 1, there exists constant $C>0$ such that 
\begin{align}
|F_n(\btha_n(\g,\h),\betas,\Sigmab) - F_n(\vb_n,\betas,\Sigmab)| \leq C \cdot\eps,\nn
\end{align}
as desired.


%
%
%
% By $f\in\rm{PL}(k)$, we have that
%\begin{align}
%|f(\sqrt{p}\btha_{n,i},\sqrt{p}\betas_i,\Sigmab_{i,i}) - f(\sqrt{p}\vb_{n,i},\sqrt{p}\betas_i,\Sigmab_{i,i}) | \leq L ( 1 +  \|\ab_i\|_2^{k-1} +  \|\bb_i\|_2^{k-1} ) \sqrt{p} |\btha_{n,i}-\vb_{n,i}|. \nn
%\end{align}•
%Thus, using Cauchy-Swchartz,
%\begin{align}
%|F_n(\sqrt{p}\btha_n,\sqrt{p}\betas,\Sigmab) - F_n(\sqrt{p}\vb_n,\sqrt{p}\betas,\Sigmab) | &\leq \frac{L}{p} \sum_{i=1}^p ( 1 +  \|\ab_i\|_2^{k-1} +  \|\bb_i\|_2^{k-1} ) \sqrt{p} |\btha_{n,i}-\vb_{n,i}|. \\
%\end{align}•
%
% Using $f\in\rm{PL}(2)$, this implies the following
%\begin{align}
%|\frac{1}{p}\sum_{i=1}^{p}f(\sqrt{p}\w_{n,i})-\frac{1}{p}\sum_{i=1}^{p}f(\sqrt{p}\vb_{n,i})| &\leq \frac{1}{p}\sum_{i=1}^{p} L(1+\sqrt{p}|\w_{n,i}|+\sqrt{p}|\vb_{n,i}|)\sqrt{p}|\w_{n,i}-\vb_{n,i}| \nn
%\\
%&\leq \nn
%\frac{L}{p}\sum_{i=1}^{p}(1+\sqrt{p}|\w_{n,i}|+\sqrt{p}|\vb_{n,i}|) \left(|\nub_i| c_1 + |\h_i|c_2\right)\nn\\
%&\leq C\left(\frac{\tn{\nub}}{\sqrt{p}}+ \frac{\tn{\h}}{\sqrt{p}}\right)(1+B+\tn{\vb_{n}}) \label{eq:fwivi}
%%&\leq C(1+2B)\left(\frac{\tn{\nub}}{\sqrt{p}}+ \frac{\tn{\h}}{\sqrt{p}}\right) 
%\end{align}\som{$\eps$ missing}
%The inequality in the first line follows from $f\in\rm{PL}(2)$. The inequality in the second line uses \eqref{eq:wivi}, and we have also recalled the notation $\nub_i=\sqrt{p}\betas_i/\sqrt{\lai}$ in Assumption \ref{ass:mu} and $\h_i=\sqrt{p}\hba_i$. The inequality in the third line uses Cauchy-Swchartz, $\tn{\w_{n}}\leq B$ and $C:=C(L,c_1,c_2)$ is a universal constant.
%
%By Assumption \ref{ass:mu}, $\tn{\nub}/{\sqrt{p}}\rP\sqrt{\E[N^2]}<\infty$. Also, with probability approaching 1: $\h/{\sqrt{p}}<2$. Similarly, it can be argued that with probability approaching 1 $\tn{\wha_n}$ is bounded by a large enough constant. Combined with \eqref{eq:fwivi}, we have shown that there exists universal constant $C'(\eps)>0$ such that with probability approaching 1:
%\begin{align}
%|\frac{1}{p}\sum_{i=1}^{p}f(\sqrt{p}\w_{n,i})-\frac{1}{p}\sum_{i=1}^{p}f(\sqrt{p}\vb_{n,i})| \leq C'(\eps) .
%\end{align}


\vspace{2pt}
\noindent\underline{Proof of \eqref{eq:ivstep2}.}~~Next, we will use Assumption \ref{ass:mu} to show that 
\begin{align}
|F_n(\vb_n,\betas,\Sigmab)| \rP \alpha_*.\label{eq:AO_conv}
\end{align}
Notice that $\vb_n$ is a function of $\betas,\Sigmab,\hba$. Concretely, define ${\ggt}:\R^3\rightarrow\R$, such that
$$
\ggt(x_1,x_2,x_3) := \frac{x_2^{-1/2}}{1+(\ksi_* x_2)^{-1}}\sqrt{\kappa}\tau_* x_3 + (1-(1+\ksi_*x_2)^{-1})x_1,
$$
and notice that
$$
\sqrt{p}\vb_{n,i} = \ggt\left(\sqrt{p}\betas_i,\lai,\h_i\right) = \frac{{\Sigmab_{i,i}^{-1/2}}}{1+(\ksi_*\Sigmab_{i,i})^{-1}}\sqrt{\kappa}\tau_*\h_i +\left(1-\frac{1}{1+\ksi_*\Sigmab_{i,i}}\right)\sqrt{p}\betas_i.
$$
Thus,
$$
F_n(\vb_n,\betas,\Sigmab) = \frac{1}{p}\sum_{i=1}^p f\left(g\left(\sqrt{p}\betas_i,\lai,\h_i\right),\sqrt{p}\betas_i,\lai\right) =: \frac{1}{p}\sum_{i=1}^p h\left(\h_i,\sqrt{p}\betas_i,\lai\right),
$$
where we have defined $h:\R^3\rightarrow\R$:
\begin{align}
h(x_1,x_2,x_3) := f\left(\ggt(x_2,x_3,x_1),x_2,x_3\right).\label{h func}
\end{align}
\cts{We will prove that $h\in\rm{PL}(4)$. Indeed, if that were the case, then Assumption \ref{ass:mu} gives}
\begin{align}
\frac{1}{p}\sum_{i=1}^p h\left(\h_i,\sqrt{p}\betas_i,\lai\right) \rP \E_{\Nn(0,1)\otimes \mu}\left[h(H,B,\Lambda)\right] &=  \E_{\Nn(0,1)\otimes \mu}\left[f\left(\ggt(B,\Lambda,H),B,\Lambda\right))\right]\\
& = \E[f\left(X_{\kappa,\sigma^2}(\Lambda,B,H),B,\Lambda\right)] = \alpha_*,
\end{align}
where the penultimate equality follows by recognizing that (cf. Eqn. \eqref{eq:X})
$$
\ggt(B,\Lambda,H) = (1-(1+\ksi_*\Lambda)^{-1})B + \sqrt{\kappa}\frac{\tau_*\Lambda^{-1/2}}{1+(\ksi_*\Lambda)^{-1}}H =  X_{\kappa,\sigma^2}(\Lambda,B,H).
$$
It remains to show that $h\in\rm{PL}(4)$. Lemma \ref{lem:hPL} in Section \ref{SM useful fact} shows that if $f\in\rm{PL}(k)$, then $h\in\rm{PL}(k+1)$ for all integers $k\geq 2$. \fy{Using this and the fact that $\Plt\subset\rm{PL}(3)$, for any $f\in\Plt$, we find that $h\in \rm{PL}(4)$. This completes the proof of \eqref{eq:ivstep2}.}
% First, consider the case $f\in\rm{PL}(2)$. Then,  $h\in\rm{PL}(3)$; thus, also $h\in\rm{PL}(4)$. Second, consider the case $f=\fl\in\Fl$. Then, we prove in Lemma \ref{lem:fL}  that $\fl\in\rm{PL}(3)$ i.e.~$\Fl\subset\rm{PL}(3)$. Thus, the desired holds in this case too.  

%In the remaining of the proof, \cts{we show that $h\in\rm{PL}(k)$} which is formalized below.

\vspace{5pt}
\subsubsection{Proof of (v):} Let $\psi:\R\rightarrow\R$ be any bounded Lipschitz function. The function $f(a,b,c) = \psi(a)$ is trivially $\rm{PL}$(2). Thus, by directly applying statement (iv) of the lemma, we find that
$$
\frac{1}{p}\sum_{i=1}^p{\psi(\sqrt{p}\btha_{n,i}(\g,\h))} \rP \E\left[\psi(X_{\kappa,\sigma^2})\right].
$$
Since this holds for any bounded Lipschitz function, we have shown that the empirical distribution of $\btha_n$ converges weakly to the distribution of $X_{\kappa,\sigma^2}$. It remains to prove boundedness of the $2$nd moment as advertised in \eqref{eq:k_AO}. Recall  from
\eqref{eq:betan} that
\begin{align}
\sqrt{p}\btha_{n,i}&=\frac{{\Sigmab_{i,i}^{-1/2}}}{1+(\ksi_n\Sigmab_{i,i})^{-1}}\sqrt{\kappa}\tau_n\h_i +\left(1-\frac{1}{1+\ksi_n\Sigmab_{i,i}}\right)(\sqrt{p}\betas_i).\nn
\end{align}
Using this, boundedness of $\Sigmab_{i,i}$ from Assumption \ref{ass:mu}, and the fact that $\tau_n\rP\tau_\star, \ksi_n\rP\ksi_\star$, there exists constant $C=C(\Sigma_{\max},\Sigma_{\min},k,\tau_\star,\ksi_\star)$ such that wpa 1,
\begin{align}
\frac{1}{p}\sum_{i=1}^p|\sqrt{p}\btha_{n,i}|^{2} \leq C\left(\frac{1}{p}\sum_{i=1}^p|\h_{i}|^{2}+\frac{1}{p}\sum_{i=1}^p|\sqrt{p}\betas_{i}|^{2}\right).\nn
\end{align}
But the two summands in the expression above are finite in the limit of $p\rightarrow\infty$. Specifically, (i) from Assumption \ref{ass:mu}, $\frac{1}{p}\sum_{i=1}^p|\sqrt{p}\betas_{i}|^{2}\rP\E[B^{2}]<\infty$; (ii) $\frac{1}{p}\sum_{i=1}^p|\h_{i}|^{2}\rP\E[H^{2}]=1$, using the facts that $\h_i\simiid\Nn(0,1)$ and $H\sim\Nn(0,1)$. This proves \eqref{eq:k_AO}, as desired.




















